\chapter*{Resumen}
\addcontentsline{toc}{chapter}{Resumen}

Las cooperativas pesqueras de Baja California necesitan anticipar el volumen de captura de
la siguiente temporada para dimensionar insumos, tripulación y compromisos de venta. Esta
tesis construye esa herramienta para la langosta roja (\emph{Panulirus interruptus}) y la
extiende a abulón y erizo, a partir de los arribos pesqueros registrados por Comunidad y
Biodiversidad, A.C.\ (COBI) y por CONAPESCA, y de variables oceanográficas satelitales
(temperatura superficial del mar y color del océano) del servicio Copernicus Marine.

El trabajo parte de una comparación de modelos de pronóstico ---ARIMA, Prophet, XGBoost,
LGBM, LSTM y su ensamble--- sobre una sola serie (langosta en San Quintín), y muestra que
ninguno explicaba la caída atípica de la temporada 2021-2022. La tesis identifica el
mecanismo faltante: las olas de calor marinas. Se implementa un índice de olas de calor
según la metodología de Hobday et al.\ (2016), que reproduce el evento ``Blob'' de 2014-2016
y el régimen cálido de 2019-2021, y se incorpora como covariable.

Sobre esa base se desarrollan tres aportaciones. Primero, un \textbf{modelo global} que
agrupa las series de especie por \ue{} del conjunto (de 28 a 33 con datos suficientes, según
el corte), con la variable objetivo en escala logarítmica, que
gana o empata contra los modelos específicos en cuatro de cinco series y rescata a las de
menor escala. Segundo, un \textbf{pronóstico probabilístico} mediante regresión cuantílica
conformalizada, que entrega intervalos con cobertura empírica cercana a la nominal y estable
incluso durante olas de calor: para una cooperativa, un rango calibrado es más accionable y
más honesto que un número puntual. Tercero, una \textbf{prueba directa de la hipótesis de
fondo} ---que el factor limitante es el volumen de datos y no la sofisticación del
modelo---: al pasar de 2 a 21 series de langosta y desbloquear las temporadas 2022-2026, la
calibración de la serie insignia se vuelve estable e invariante a la composición del pool,
mientras que ni la selección de variables por SHAP, ni la optimización bayesiana, ni un
Temporal Fusion Transformer lo consiguen. El resultado se materializa en una aplicación de
inferencia que sirve el intervalo calibrado por especie y \ue{}.

\vspace{0.5cm}
\noindent\textbf{Palabras clave:} pronóstico de series de tiempo, pesquerías, langosta roja,
Baja California, olas de calor marinas, predicción conformal, modelos globales,
interpretabilidad.

\cleardoublepage

\chapter*{Abstract}
\addcontentsline{toc}{chapter}{Abstract}

Small-scale fishing cooperatives in Baja California need to anticipate next season's catch
volume in order to plan inputs, crew and sales commitments. This thesis builds that tool for
red spiny lobster (\emph{Panulirus interruptus}), extending it to abalone and sea urchin,
using landings recorded by Comunidad y Biodiversidad, A.C.\ (COBI) and CONAPESCA together
with satellite oceanographic covariates (sea surface temperature and ocean colour) from
Copernicus Marine.

Starting from a comparison of forecasting models ---ARIMA, Prophet, XGBoost, LGBM, LSTM and
their ensemble--- on a single series, we show that none of them explained the anomalous
2021-2022 season. We identify the missing mechanism as marine heatwaves, implement the
Hobday et al.\ (2016) index ---which recovers the 2014-2016 ``Blob'' and the 2019-2021 warm
regime--- and add it as a covariate.

Three contributions follow. A \textbf{global model} pooling the dataset's species-by-cooperative series (28--33 with
enough data, depending on the cut) with a log-scaled target matches or beats series-specific models in four of five series and
rescues the small-scale ones. A \textbf{probabilistic forecast} based on conformalized
quantile regression yields intervals whose empirical coverage stays close to nominal even
during marine heatwaves; for a cooperative, a calibrated range is more actionable ---and
more honest--- than a point estimate. Finally, a \textbf{direct test of the underlying
hypothesis} ---that the binding constraint is data volume rather than model
sophistication---: expanding from 2 to 21 lobster series and unlocking the 2022-2026 seasons
makes the flagship series' calibration stable and invariant to pool composition, whereas
SHAP-based feature selection, Bayesian hyperparameter search and a Temporal Fusion
Transformer all fail to do so. The result is delivered as an inference application serving the
calibrated interval per species and cooperative.

\vspace{0.5cm}
\noindent\textbf{Keywords:} time series forecasting, fisheries, red spiny lobster, Baja
California, marine heatwaves, conformal prediction, global models, interpretability.

\cleardoublepage
